\documentclass[10pt,a4paper]{amsart}
\usepackage[margin=19mm]{geometry}
\usepackage{amsmath,amssymb,mathtools}
\usepackage[scr=rsfs,cal=euler]{mathalfa}
\usepackage{xfrac}
\usepackage[unicode,hidelinks,bookmarks=false]{hyperref}
\newcommand{\ie}{i.e.\ }
\newcommand{\cf}{cf.\ }
\newcommand{\resp}{respectively\ }
\newcommand{\Subsection}{\subsection}
\providecommand{\nobreakdash}{\nobreak\hbox{-}}
\setlength{\emergencystretch}{2em}
\multlinegap0pt
\usepackage{bm}
\usepackage{bbm}
\usepackage{dsfont}
\usepackage{xspace}
\usepackage{cancel}
\usepackage{mathtools}
\usepackage{amsthm}
\makeatletter
\@ifundefined{theorem}{\newtheorem{theorem}{Theorem}}{}
\@ifundefined{lemma}{\newtheorem{lemma}[theorem]{Lemma}}{}
\makeatother
\newcommand{\C}{\mathbb{C}}
\newcommand{\R}{\mathbb{R}}
\newcommand{\cA}{\mathcal{A}}
\newcommand{\rd}{\mathrm{d}}
\newcommand{\ri}{\mathrm{i}}
\title[Integral equation methods for scattering by general compact ]{Integral equation methods for scattering by general compact obstacles: wavenumber-explicit estimates}
\author{Simon N. Chandler-Wilde, Siavash Sadeghi}
\date{Source: arXiv:2601.19456v2 (CC BY 4.0)}
\begin{document}
\maketitle

\noindent\emph{Source and licence.} Integral equation methods for scattering by general compact obstacles: wavenumber-explicit estimates. Version arXiv:2601.19456v2 (CC BY 4.0), distributed under the Creative Commons Attribution 4.0 International licence, as recorded in the arXiv licence metadata for this exact version and shown on the abstract page https://arxiv.org/abs/2601.19456v2. Full rights notice: licenses/arxiv-2601.19456-CC-BY-4.0.txt.\par\smallskip

\noindent\textit{Editorial scope.} Counted unit: the Lemma whose source label is \texttt{lem:step2} in the article (source0.tex lines 955--961, source label \texttt{lem:step2}), delivered together with the article's own complete original proof of it (source0.tex lines 962--992, one \texttt{proof} environment of 31 lines). No mathematics is written, repaired or re-derived: statement and proof are the article's own text line for line. Required context retained uncounted: lem:step1 (lines 931--936). Every definition, display and label of those lines is retained unchanged. Omitted from this package and not redistributed: 1--930, 937--954, 993--1745. Nothing inside a retained excerpt is omitted or altered: every retained body is delivered exactly as the article's lines are written, so no omission receipt of any kind accompanies this package. Citations: the retained excerpts carry no citation command at all, so no bibliography entry of the article is transcribed and no reference is redistributed. Reference adaptation: none was needed and none was made; every reference inside the delivered unit resolves inside it, so no unresolved reference or citation is produced. Printed slips kept literal: no printed word, symbol, hypothesis or number of the counted statement or of its proof was altered. Layout and class: the article's own class is \texttt{birkjour} with packages including graphicx, xcolor, soul, amsmath, amsthm, amssymb, tikz, hyperref; the wrapper is the lane's pinned \texttt{amsart} editorial wrapper at 10pt on A4, which restates the theorem-like environments the retained text needs and adds the article's own macro definitions that the retained text needs (\texttt{\textbackslash C}, \texttt{\textbackslash R}, \texttt{\textbackslash cA}, \texttt{\textbackslash rd}, \texttt{\textbackslash ri}), with extra wrapper packages (bm, bbm, dsfont, xspace, cancel, mathtools). The delivered numbering is the unit's own. Assets: no retained span contains a figure environment, a graphics-inclusion command, a PDF-page inclusion, a table or a TikZ picture; no image file is copied into this package, the package declares no asset dependency and no image file is redistributed with it. Licence: version arXiv:2601.19456v2 of this article is distributed under the Creative Commons Attribution 4.0 International licence, as recorded in the arXiv licence metadata for this exact version and shown on the abstract page https://arxiv.org/abs/2601.19456v2. A full rights notice is delivered at licenses/arxiv-2601.19456-CC-BY-4.0.txt. Characters: every retained excerpt is pure ASCII, so the delivered engine, XeLaTeX with its default Latin Modern text font, reports no missing character.
\begin{lemma} \label{lem:step1} Given any $k_0>0$ and $R>R_\Gamma$ there exists $C_1>0$ such that
$$
\|\phi\|_{H_k^{-1}(\R^n)} \leq C_1 \|u\|_{H^1_k(B_R)},
$$
if $\phi\in H^{-1}_\Gamma$, $k\geq k_0$, and $u:= \cA_\lambda \phi$, where $\lambda\in \C$ with $\mathrm{Im}\,\lambda \geq 0$ and $\lambda^2=k^2+\ri \beta(k)$, with $0\leq \beta(k)\leq k^2$. 
\end{lemma}

\begin{lemma} \label{lem:step2} Suppose that  $k_0>0$. Then  there exists $C_2>0$ such that, for all  $k\geq k_0$, $g\in \widetilde H^1(\Gamma^c)^{\perp_k}$, and 
$\lambda\in \C$ with $\lambda_i := \mathrm{Im}\, \lambda >0$ and $\lambda^2=k^2+\ri \beta(k)$, with $0<\beta(k)\leq k^2$,  it holds that
$$
\|w\|_{H^1_k(\R^n)} \leq \frac{C_2k^2}{\beta(k)} \|g\|_{H^1_k(\R^n)},
$$ 
where $w:= \cA_\lambda \mu$ and $\mu:= A_{k,\lambda}^{-1}g\in H^{-1}_\Gamma$.
\end{lemma}

\begin{proof} Suppose that $k\geq k_0$ and $\lambda$, $g$, $\mu$, and $w$ are as in the statement of the lemma. Let us define $b_\lambda: H_k^1(\mathbb{R}^n)\times H_k^1(\mathbb{R}^n)\rightarrow\mathbb{C}$ by
\begin{equation*}
 b_\lambda(u,v):=\int_{\R^n}\left(\lambda^2 u \bar v - \nabla u\cdot \nabla \bar v\right)\, \rd x, \qquad u,v \in H^1(\R^n).
\end{equation*}
This is the sesquilinear form that appears in the weak form of the Helmholtz equation on $\R^n$: 
if $u\in H^1(\R^n)$ and $f\in H^{-1}(\R^n)$, then $\Delta u + \lambda^2 u = -f$ in a distributional sense if and only if
$$
b_\lambda(u,v) = -\langle f,v\rangle, \qquad \forall v\in H^1(\R^n).
$$
In particular, since $(\Delta + \lambda^2)w = - \mu$ and $\langle \mu, v\rangle =0$ for all $v\in \widetilde H^1(\Gamma^c)$, we have that
\begin{align*}
|b_\lambda(w,w)| & = |\langle \mu,\cA_\lambda \mu\rangle| = |\langle \mu, A_{k,\lambda} \mu\rangle|=|\langle \mu,g\rangle|\\ & \leq \|\mu\|_{H_k^{-1}(\R^n)} \|g\|_{H^1_k(\R^n)} \leq C_1\|w\|_{H_k^{1}(\R^n)} \|g\|_{H^1_k(\R^n)},
\end{align*}
by Lemma \ref{lem:step1}. Further, for $u\in H^1(\R^n)$,
$$
\bar \lambda b_\lambda(u,u) = \int_{\R^n}\left(\lambda |\lambda|^2 |u|^2  - \bar\lambda |\nabla u|^2\right)\, \rd x
$$
so that
$$
|\bar \lambda b_\lambda(u,u)| \geq \mathrm{Im}( \bar \lambda b_\lambda(u,u)) = \lambda_i \int_{\R^n}\left(|\lambda|^2 |u|^2  + |\nabla u|^2\right)\, \rd x \geq \lambda_i \|u\|_{H^1_k(\R^n)}^2,
$$
i.e., $b_\lambda(\cdot,\cdot)$ is coercive, and
$$
\lambda_i \|w\|_{H_k^{1}(\R^n)} \leq C_1|\lambda| \|g\|_{H^1_k(\R^n)}.
$$
Further, $\beta(k)\leq k^2$ so that $|\lambda| \leq \sqrt{2}\,k$ and
$$
\lambda_i = \frac{\beta(k)}{\sqrt{2}\sqrt{k^2+\sqrt{k^4+\beta^2(k)}}}\geq \frac{\beta(k)}{\sqrt{2}\sqrt{1+\sqrt{2}}\,k},
$$
from which the claimed result follows.
\end{proof}

\end{document}
